\documentclass[11pt]{article}
\usepackage[utf8]{inputenc}
\usepackage[margin=1in]{geometry}
\usepackage{amsmath}
\usepackage{graphicx}
\usepackage{booktabs}
\usepackage{hyperref}
\usepackage{natbib}
\usepackage{lineno}
\usepackage{xcolor}
\usepackage{multirow}
\usepackage{float}

\linenumbers

\title{Profiling of Chimeric RNAs in Human Retinal Development with Retinal Organoids}

\author{
Author~One$^{1,\dagger}$, Author~Two$^{1,\dagger}$, Author~Three$^{2}$, Author~Four$^{1}$, \\
Author~Five$^{2}$, Author~Six$^{1,*}$ \\[6pt]
$^{1}$Department of Ophthalmology, Institute of Visual Science, China \\
$^{2}$Department of Molecular Biology, Institute of Genetics, China \\[4pt]
$^{\dagger}$These authors contributed equally. \\
$^{*}$Corresponding author: corresponding@example.edu
}

\date{}

\begin{document}
\maketitle

% ============================================================
\begin{abstract}
Chimeric RNAs have emerged as important regulators of stem cell differentiation and central nervous system development, yet their landscape during human retinal development remains uncharted. Here we present the first comprehensive expression atlas of chimeric RNAs across human retinal organoid (RO) development, spanning 22 samples at eight developmental stages from day 0 (D0) to D120. Using FusionCatcher-based detection with stringent filtering, we identified hundreds of chimeric RNA events and found that intra-chromosomal chimeric RNAs increased significantly with developmental progression ($r^{2} = 0.93$, $p = 7.6 \times 10^{-4}$, Pearson correlation), while inter-chromosomal events showed no such trend. Protein-coding gene pairs constituted 61.9\% of chimeric RNA parental genes, and in-frame chimeric RNAs accounted for 11.2\% of all events. We validated four isoforms of the chimeric RNA CTNNBIP1-CLSTN1 (CTCL) by PCR and Sanger sequencing in D60 organoids. Loss-of-function experiments revealed that CTCL knockdown obstructed neural retinal differentiation while promoting retinal pigment epithelium (RPE) fate, with thinner outer layers at D60. Transcriptomic analysis of CTCL-knockdown organoids unexpectedly revealed significant alterations in epithelial-mesenchymal transition, extracellular matrix organization, and epithelial cell differentiation pathways rather than the anticipated Wnt signaling changes. These findings establish chimeric RNAs as previously unrecognized regulators of retinogenesis and identify CTCL as a key modulator of the neural retina versus RPE fate decision.

\vspace{0.5em}
\noindent\textbf{Keywords:} chimeric RNA, retinal organoid, retinal development, CTCL, RPE differentiation, epithelial-mesenchymal transition
\end{abstract}

% ============================================================
\section{Introduction}

Chimeric RNAs---transcripts formed by the fusion of sequences from two or more distinct genomic loci---have attracted growing attention as regulators of cellular differentiation and tissue development \citep{li2019chimeric, wu2019circular}. Although initially discovered in cancer contexts, mounting evidence demonstrates that chimeric RNAs are expressed in normal tissues and play physiological roles during development \citep{tang2017fusion}. In particular, chimeric RNAs have been implicated in neural stem and progenitor cell differentiation within the central nervous system (CNS) \citep{li2017chimeric_cns, gao2022ctcl}.

The retina, as a developmentally and anatomically integral component of the CNS, undergoes a precisely orchestrated differentiation program in which retinal progenitor cells give rise to seven major cell types in a conserved temporal order \citep{cepko2014retinal}. Despite substantial progress in understanding the transcriptomic landscape of retinal development at the single-cell level \citep{cowan2020cell, macosko2015highly}, the chimeric RNA transcriptome of the developing retina has not been explored. This gap is significant because chimeric RNAs may regulate the balance between proliferation and differentiation in retinal progenitors, analogous to their roles in cerebral cortex development.

Retinal organoids (ROs) derived from human embryonic stem cells (hESCs) provide an accessible model system that recapitulates the spatiotemporal features of human retinogenesis \citep{nakano2012retinal, zhong2014generation, reichman2017generation}. These three-dimensional structures develop laminated retinal tissue over a period of weeks to months, generating photoreceptors, ganglion cells, and retinal pigment epithelium (RPE) in a manner that closely parallels in vivo development \citep{yue2021retinal_organoid, kallman2020retinal}.

Among chimeric RNAs with demonstrated developmental functions, CTNNBIP1-CLSTN1 (CTCL) is of particular interest. CTCL was recently shown to regulate cerebral cortex development, where its downregulation causes premature neuronal differentiation in cerebral organoids \citep{gao2022ctcl}. The upstream parental gene, CTNNBIP1, encodes an inhibitor of the Wnt/\(\beta\)-catenin pathway \citep{zhou2016wnt}, while the downstream gene, CLSTN1, encodes a transmembrane protein involved in vesicular trafficking \citep{adijanto2014src}. Whether CTCL functions in retinogenesis---another CNS developmental process---has not been tested.

In this study, we present the first expression atlas of chimeric RNAs throughout RO development from D0 to D120 across 22 samples at eight developmental stages. We characterize the genomic distribution, predicted functional effects, and parental gene biotypes of chimeric RNAs at each stage. We identify and validate four CTCL isoforms in retinal organoids and demonstrate through loss-of-function experiments that CTCL knockdown obstructs neural retinal differentiation while promoting RPE fate. Unexpectedly, transcriptomic analysis reveals that CTCL knockdown alters epithelial-mesenchymal transition (EMT) and extracellular matrix pathways rather than the Wnt signaling pathway predicted by the function of its parental gene CTNNBIP1.

% ============================================================
\section{Results}

\subsection{Chimeric RNA landscape during retinal organoid development}

To profile chimeric RNAs across retinal development, we performed RNA sequencing on 22 RO samples spanning eight developmental stages (D0, D30, D45, D60, D75, D90, D105, D120) and analyzed the data using FusionCatcher with a stringent filter requiring at least one spanning unique read \citep{nicorici2014fusioncatcher}. Expression levels were quantified as $\log(\text{spanning unique reads} / \text{total reads} \times 10^{7})$.

Chimeric RNAs were continuously expressed throughout all developmental stages. Analysis of genomic distribution revealed a striking dichotomy between intra- and inter-chromosomal events (Table~\ref{tab:chimeric_landscape}).

\begin{table}[H]
\centering
\caption{Chimeric RNA counts by genomic category across developmental stages. Pearson correlation with developmental stage (D0--D120) is shown. $^{**}p < 0.001$.}
\label{tab:chimeric_landscape}
\begin{tabular}{lcccccccc}
\toprule
\textbf{Category} & \textbf{D0} & \textbf{D30} & \textbf{D45} & \textbf{D60} & \textbf{D75} & \textbf{D90} & \textbf{D105} & \textbf{D120} \\
\midrule
Intra-chromosomal & 42 & 58 & 71 & 89 & 102 & 118 & 131 & 147 \\
Inter-chromosomal & 31 & 28 & 35 & 33 & 29 & 37 & 32 & 34 \\
Total events & 73 & 86 & 106 & 122 & 131 & 155 & 163 & 181 \\
\midrule
\multicolumn{9}{l}{\textit{Intra-chromosomal:} $r^{2} = 0.93$, $p = 7.6 \times 10^{-4}$ $^{**}$ (Pearson)} \\
\multicolumn{9}{l}{\textit{Inter-chromosomal:} $r^{2} = 0.04$, $p = 0.63$ (Pearson)} \\
\bottomrule
\end{tabular}
\end{table}

Intra-chromosomal chimeric RNAs increased significantly with RO developmental progression ($r^{2} = 0.93$, $p = 7.6 \times 10^{-4}$, Pearson correlation), while inter-chromosomal chimeric RNAs showed no significant trend ($r^{2} = 0.04$, $p = 0.63$).

\subsection{Characterization of chimeric RNA types and parental gene biotypes}

We next classified all detected chimeric RNAs by their predicted functional effect and parental gene biotypes (Table~\ref{tab:chimeric_types}).

\begin{table}[H]
\centering
\caption{Distribution of chimeric RNAs by predicted functional effect across all 22 samples. Proportions are computed relative to total events. Chi-squared goodness-of-fit test: $\chi^{2} = 284.3$, $df = 13$, $p < 1 \times 10^{-10}$.}
\label{tab:chimeric_types}
\begin{tabular}{lcc}
\toprule
\textbf{Predicted Effect Category} & \textbf{Count} & \textbf{Percentage (\%)} \\
\midrule
In frame & 94 & 11.2 \\
CDS(truncated)\_UTR & 90 & 10.7 \\
UTR\_UTR & 90 & 10.7 \\
CDS\_CDS (out of frame) & 78 & 9.3 \\
UTR\_CDS(truncated) & 72 & 8.6 \\
CDS(truncated)\_CDS & 65 & 7.8 \\
Intronic\_CDS & 58 & 6.9 \\
CDS\_Intronic & 52 & 6.2 \\
Intronic\_Intronic & 47 & 5.6 \\
UTR\_Intronic & 41 & 4.9 \\
Intronic\_UTR & 38 & 4.5 \\
CDS\_UTR & 34 & 4.1 \\
UTR\_CDS & 30 & 3.6 \\
Other & 49 & 5.9 \\
\midrule
\textbf{Total} & \textbf{838} & \textbf{100.0} \\
\bottomrule
\end{tabular}
\end{table}

In-frame chimeric RNAs, which potentially encode novel fusion proteins, accounted for 11.2\% of events. The majority (61.9\%) of chimeric RNAs were formed between two protein-coding genes, with additional contributions from lncRNA, pseudogene, and other non-coding gene combinations (Table~\ref{tab:biotype}).

\begin{table}[H]
\centering
\caption{Parental gene biotype combinations of chimeric RNAs. Fisher's exact test comparing protein-coding pairs vs.\ all other combinations: OR $= 3.41$, 95\% CI [2.78, 4.19], $p < 1 \times 10^{-15}$.}
\label{tab:biotype}
\begin{tabular}{lcc}
\toprule
\textbf{Biotype Combination} & \textbf{Count} & \textbf{Percentage (\%)} \\
\midrule
Protein-coding + Protein-coding & 519 & 61.9 \\
Protein-coding + lncRNA & 108 & 12.9 \\
lncRNA + Protein-coding & 72 & 8.6 \\
Protein-coding + Pseudogene & 48 & 5.7 \\
lncRNA + lncRNA & 35 & 4.2 \\
Pseudogene + Protein-coding & 28 & 3.3 \\
Other combinations & 28 & 3.4 \\
\midrule
\textbf{Total} & \textbf{838} & \textbf{100.0} \\
\bottomrule
\end{tabular}
\end{table}

Analysis of sequence motifs around fusion sites using the seqLogo R package \citep{schneider2010seqlogo} revealed no strong consensus sequence. Spearman correlation between chimeric RNA expression and parental gene expression (mapped with Hisat2 and quantified with featureCounts \citep{kim2019hisat2, liao2014featurecounts}) was not significant ($\rho = 0.08$, $p = 0.31$), indicating that chimeric RNA expression is regulated independently of parental gene transcription.

\subsection{Identification and validation of CTCL isoforms in retinal organoids}

Given the established role of CTCL in cerebral development \citep{gao2022ctcl}, we investigated its expression across RO development. We detected four distinct CTCL isoforms by RNA-seq spanning reads and validated all four by PCR and Sanger sequencing in D60 organoids (Table~\ref{tab:ctcl_isoforms}).

\begin{table}[H]
\centering
\caption{CTCL isoform expression across retinal organoid development. Expression is quantified as $\log(\text{spanning unique reads}/\text{total reads} \times 10^{7})$. Two-way ANOVA (isoform $\times$ stage): isoform effect $F_{3,56} = 18.7$, $p = 2.1 \times 10^{-8}$; stage effect $F_{7,56} = 4.2$, $p = 8.4 \times 10^{-4}$.}
\label{tab:ctcl_isoforms}
\begin{tabular}{lcccccccc}
\toprule
\textbf{Isoform} & \textbf{D0} & \textbf{D30} & \textbf{D45} & \textbf{D60} & \textbf{D75} & \textbf{D90} & \textbf{D105} & \textbf{D120} \\
\midrule
In-frame & 1.82 & 2.14 & 2.47 & \textbf{2.91} & 2.63 & 2.38 & 2.11 & 1.94 \\
UTR\_CDS(trunc)-1 & 1.21 & 1.54 & 1.78 & 2.03 & 1.89 & 1.65 & 1.42 & 1.31 \\
CDS(trunc)\_Intronic & 0.87 & 1.02 & 1.19 & 1.38 & 1.27 & 1.14 & 0.98 & 0.91 \\
UTR\_CDS(trunc)-2 & 0.64 & 0.78 & 0.93 & 1.11 & 1.02 & 0.89 & 0.76 & 0.68 \\
\bottomrule
\end{tabular}
\end{table}

The in-frame isoform showed peak expression at D60, coinciding with the critical window of photoreceptor precursor specification. All four isoforms were confirmed by PCR amplification and Sanger sequencing of the fusion junctions.

\subsection{CTCL knockdown disrupts retinal organoid differentiation}

To test the functional role of CTCL in retinogenesis, we used lentiviral shRNA to knock down CTCL in CRX-tdTomato reporter hES-derived ROs at D12, collecting organoids at D60 for analysis ($n = 3$ independent experiments). Knockdown efficiency was confirmed at 48--72 hours post-infection. Importantly, CTCL knockdown did not alter expression of either parental gene (CTNNBIP1: $p = 0.71$; CLSTN1: $p = 0.64$; two-tailed $t$-test, $n = 4$), confirming the specificity of the shRNA for the chimeric transcript.

Phenotypic analysis revealed that CTCL-knockdown organoids had thinner outer layers at D60 compared to scramble controls, with altered CRX-tdTomato signal distribution (Table~\ref{tab:kd_phenotype}).

\begin{table}[H]
\centering
\caption{Immunohistochemical quantification of marker expression in CTCL-knockdown vs.\ scramble control organoids at D60. Fluorescence intensity is normalized to DAPI. Two-tailed unpaired $t$-test, $n = 3$ independent experiments per group.}
\label{tab:kd_phenotype}
\begin{tabular}{lccccc}
\toprule
\textbf{Marker} & \textbf{Category} & \textbf{shCTCL (mean $\pm$ SD)} & \textbf{Scramble (mean $\pm$ SD)} & \textbf{$t$-statistic} & \textbf{$p$-value} \\
\midrule
CRX & Photoreceptor precursor & $0.42 \pm 0.08$ & $0.78 \pm 0.11$ & $-4.58$ & $0.010$ \\
OTX2 & Neural/retinal & $0.51 \pm 0.09$ & $0.84 \pm 0.12$ & $-3.82$ & $0.019$ \\
VSX2 & Retinal progenitor & $0.63 \pm 0.10$ & $0.91 \pm 0.07$ & $-3.96$ & $0.017$ \\
RCVRN & Photoreceptor & $0.38 \pm 0.07$ & $0.69 \pm 0.09$ & $-4.72$ & $0.009$ \\
MITF & RPE-specific & $0.89 \pm 0.14$ & $0.31 \pm 0.06$ & $6.62$ & $0.003$ \\
RPE65 & RPE-specific & $0.76 \pm 0.11$ & $0.22 \pm 0.05$ & $7.71$ & $0.002$ \\
PAX6 & Eye field/progenitor & $0.71 \pm 0.13$ & $0.54 \pm 0.10$ & $1.80$ & $0.146$ \\
\bottomrule
\end{tabular}
\end{table}

CTCL knockdown significantly reduced photoreceptor and neural retinal markers (CRX, OTX2, VSX2, RCVRN; all $p < 0.02$) while dramatically increasing RPE-specific markers (MITF, RPE65; both $p < 0.005$). PAX6, which marks both retinal progenitors and RPE, showed a non-significant trend toward increase.

\subsection{Transcriptomic analysis reveals unexpected pathway alterations}

RNA sequencing of CTCL-knockdown organoids at D60 identified differentially expressed genes (DEGs) defined as $|\log_{2}(\text{shCTCL}/\text{scramble})| > 1$ with $p < 0.05$ (Table~\ref{tab:degs}). Gene set enrichment analysis (GSEA) was performed using clusterProfiler \citep{wu2021clusterprofiler} with the c8 reference gene set \citep{subramanian2005gsea}.

\begin{table}[H]
\centering
\caption{Top enriched pathways in CTCL-knockdown organoids. GSEA normalized enrichment scores (NES) and adjusted $p$-values (Benjamini--Hochberg) are shown. $\uparrow$ indicates upregulated in shCTCL; $\downarrow$ indicates downregulated.}
\label{tab:degs}
\begin{tabular}{llccc}
\toprule
\textbf{Direction} & \textbf{Pathway} & \textbf{NES} & \textbf{Adj.\ $p$} & \textbf{Gene Count} \\
\midrule
$\uparrow$ & Epithelial-mesenchymal transition & $2.41$ & $3.2 \times 10^{-6}$ & 187 \\
$\uparrow$ & Extracellular matrix organization & $2.28$ & $8.7 \times 10^{-6}$ & 154 \\
$\uparrow$ & Epithelial cell differentiation & $2.15$ & $2.1 \times 10^{-5}$ & 132 \\
$\uparrow$ & RPE cell gene set (c8) & $1.98$ & $7.4 \times 10^{-5}$ & 98 \\
$\uparrow$ & Collagen fibril organization & $1.87$ & $1.8 \times 10^{-4}$ & 67 \\
$\downarrow$ & Neurogenesis & $-2.34$ & $5.1 \times 10^{-6}$ & 201 \\
$\downarrow$ & Nervous system development & $-2.19$ & $1.4 \times 10^{-5}$ & 178 \\
$\downarrow$ & Morphogenesis of a branching structure & $-1.92$ & $5.8 \times 10^{-5}$ & 89 \\
$\downarrow$ & Axon guidance & $-1.78$ & $2.3 \times 10^{-4}$ & 74 \\
$\downarrow$ & Photoreceptor cell differentiation & $-1.65$ & $6.1 \times 10^{-4}$ & 52 \\
\midrule
\multicolumn{5}{l}{\textit{Not significantly changed:} Wnt signaling ($p = 0.42$), Vesicular transport ($p = 0.58$)} \\
\bottomrule
\end{tabular}
\end{table}

Critically, neither Wnt signaling (expected based on CTNNBIP1 function) nor vesicular transport (expected based on CLSTN1 function) was significantly altered, demonstrating that CTCL acts through mechanisms independent of its parental gene pathways.

\subsection{RPE marker gene upregulation in CTCL-knockdown organoids}

Consistent with the RPE-promoting phenotype, RNA-seq confirmed significant upregulation of key RPE transcription factors and functional genes in CTCL-knockdown organoids (Table~\ref{tab:rpe_genes}).

\begin{table}[H]
\centering
\caption{RPE marker gene expression changes upon CTCL knockdown. $\log_{2}$ fold change (shCTCL / scramble) with DESeq2-derived adjusted $p$-values \citep{love2014deseq2}. $\uparrow$ denotes upregulation in shCTCL.}
\label{tab:rpe_genes}
\begin{tabular}{lcccc}
\toprule
\textbf{Gene} & \textbf{Function} & \textbf{$\log_{2}$FC} & \textbf{Adj.\ $p$} & \textbf{Direction} \\
\midrule
MITF & RPE master TF & $2.87$ & $1.4 \times 10^{-8}$ & $\uparrow$ \\
RPE65 & Visual cycle & $3.21$ & $4.7 \times 10^{-9}$ & $\uparrow$ \\
BEST1 & Ion channel & $2.54$ & $8.3 \times 10^{-7}$ & $\uparrow$ \\
TYR & Melanogenesis & $2.93$ & $2.1 \times 10^{-8}$ & $\uparrow$ \\
TYRP1 & Melanogenesis & $2.41$ & $3.6 \times 10^{-7}$ & $\uparrow$ \\
PMEL & Melanosome biogenesis & $2.18$ & $1.2 \times 10^{-6}$ & $\uparrow$ \\
RLBP1 & Visual cycle & $1.87$ & $5.4 \times 10^{-6}$ & $\uparrow$ \\
LRAT & Visual cycle & $1.64$ & $2.8 \times 10^{-5}$ & $\uparrow$ \\
\bottomrule
\end{tabular}
\end{table}

% ============================================================
\section{Discussion}

Our study provides the first comprehensive atlas of chimeric RNAs during human retinal organoid development and identifies CTCL as a key regulator of the neural retina versus RPE fate decision. Several findings merit discussion.

The significant increase of intra-chromosomal chimeric RNAs with development ($r^{2} = 0.93$, $p = 7.6 \times 10^{-4}$) suggests that cis-splicing-mediated chimeric RNA production becomes more active as retinal cells diversify and mature. This parallels observations in other developing tissues where chromatin reorganization during differentiation may facilitate read-through transcription between adjacent genes \citep{tang2017fusion}. The absence of a similar trend for inter-chromosomal events suggests that trans-splicing mechanisms do not show the same developmental regulation.

The predominance of protein-coding parental gene pairs (61.9\%) implies functional relevance, as in-frame fusion products could generate novel proteins with combinatorial domain architectures. The 11.2\% in-frame rate we observed is consistent with rates reported in other normal tissues \citep{li2019chimeric}.

Perhaps the most striking finding is the phenotypic consequence of CTCL knockdown. While CTCL downregulation was previously shown to cause premature neuronal differentiation in cerebral organoids \citep{gao2022ctcl}, we observed a qualitatively different effect in retinal organoids: CTCL knockdown obstructed neural retinal differentiation and instead promoted RPE fate. This tissue-specific difference highlights the context-dependent nature of chimeric RNA function. The retina presents a unique developmental choice---neural retina versus RPE---that is not paralleled in the cerebral cortex, and CTCL appears to regulate this binary decision.

The unexpected activation of EMT and extracellular matrix pathways, rather than Wnt signaling, upon CTCL knockdown suggests that the chimeric RNA has acquired functions distinct from its parental genes. RPE differentiation involves EMT-related processes \citep{thomas2019emt, fuhrmann2014rpe}, and the upregulation of these pathways is consistent with the RPE-promoting phenotype. This finding illustrates a general principle: chimeric RNAs may function through mechanisms that cannot be predicted from their parental gene functions.

Our study has limitations. The use of retinal organoids, while providing a tractable model for human retinal development, does not fully recapitulate the signaling environment of the developing eye. The number of independent shRNA experiments ($n = 3$) is modest, and future studies should include rescue experiments and alternative knockdown strategies. Additionally, the mechanism by which CTCL regulates EMT pathways remains to be elucidated.

% ============================================================
\section{Methods}

\subsection{Retinal organoid generation and culture}

Human embryonic stem cells (CRX-tdTomato reporter line) were differentiated into retinal organoids following established protocols \citep{nakano2012retinal, zhong2014generation}. Organoids were cultured from D0 to D120, with samples collected at eight stages (D0, D30, D45, D60, D75, D90, D105, D120). A total of 22 organoid samples were collected across all stages, with 2--3 biological replicates per stage.

\subsection{RNA sequencing and chimeric RNA detection}

Total RNA was extracted from each sample and sequenced using the VAHTS Universal V6 RNA-seq Library Prep Kit for Illumina. Chimeric RNAs were identified using FusionCatcher software \citep{nicorici2014fusioncatcher} with default parameters and a stringency filter requiring at least one spanning unique read. Expression levels were quantified as $\log(\text{spanning unique reads} / \text{total reads} \times 10^{7})$. Reads were aligned to the hg38 reference genome using Hisat2 \citep{kim2019hisat2}, and transcript-level expression was quantified with featureCounts \citep{liao2014featurecounts}.

\subsection{PCR validation and Sanger sequencing}

CTCL isoforms were validated by RT-PCR using isoform-specific primers designed to span the fusion junctions. PCR products were gel-purified and submitted for Sanger sequencing. Sequences were aligned to the predicted fusion junctions to confirm identity.

\subsection{CTCL knockdown}

Lentiviral shRNA constructs targeting CTCL (shCTCL) and scramble controls were packaged in HEK293T cells. ROs were transduced at D12 at a multiplicity of infection of 5. Knockdown efficiency was confirmed by RT-qPCR at 48--72 hours post-infection. Three independent experiments were performed for each condition. Organoids were collected at D60 for downstream analyses.

\subsection{Immunohistochemistry}

D60 organoids were fixed in 4\% paraformaldehyde, embedded in OCT, and cryosectioned at 12 $\mu$m. Sections were stained with antibodies against CRX, OTX2, VSX2, RCVRN, MITF, RPE65, and PAX6. Fluorescence intensity was quantified using ImageJ and normalized to DAPI signal.

\subsection{Differential expression and pathway analysis}

RNA-seq data from shCTCL and scramble organoids at D60 were analyzed using DESeq2 \citep{love2014deseq2}. DEGs were defined as $|\log_{2}\text{FC}| > 1$ and adjusted $p < 0.05$. Functional enrichment of upregulated and downregulated gene sets was performed using Metascape. GSEA was conducted using clusterProfiler \citep{wu2021clusterprofiler} with the c8 cell-type signature reference gene set \citep{subramanian2005gsea}. Sequence motifs surrounding fusion sites were analyzed with the seqLogo R package \citep{schneider2010seqlogo}.

\subsection{Statistical analysis}

All statistical tests were performed in R (version 4.2). Pearson correlations were used for developmental trend analyses. Two-tailed unpaired $t$-tests were used for IHC comparisons between shCTCL and scramble groups. Two-way ANOVA was used for multi-factor analyses. Spearman correlation was used for expression relationship analyses. Multiple comparisons were corrected using the Benjamini--Hochberg method unless otherwise stated.

% ============================================================
\section{Conclusion}

We present the first chimeric RNA atlas of human retinal organoid development and demonstrate that the chimeric RNA CTCL regulates the balance between neural retina and RPE fate. The developmental increase of intra-chromosomal chimeric RNAs, the predominance of protein-coding parental gene pairs, and the tissue-specific function of CTCL collectively establish chimeric RNAs as a previously unrecognized layer of gene regulation in retinogenesis. These findings open new avenues for understanding retinal cell fate decisions and may inform strategies for directed differentiation of retinal cell types from pluripotent stem cells.

\bibliographystyle{plainnat}
\bibliography{references}

\end{document}
